Large-scale copy number alterations (CNAs) can be inferred from single-cell RNA-seq (scRNA-seq) data, which makes it possible to distinguish malignant from non-malignant cells and to define genetic subclones in the same data used to study cell states. The approach was introduced in single-cell studies of glioblastoma and was used in a series of subsequent tumour atlases. inferCNV is an R/Bioconductor package that implements and extends it, and it has been widely used since 2018. inferCNV compares the expression of genes ordered along the genome in tumour cells with that of a user-defined set of normal reference cells. It produces reference-normalised, positionally smoothed residual expression profiles; groups cells into subclusters; predicts discrete copy number states with a six-state hidden Markov model calibrated by in-silico spike-in simulation; and filters predicted CNAs with a Bayesian mixture model. Here we describe the algorithm and its reference implementation and illustrate its behaviour on simulated data and on tumours with DNA-based copy number measurements. Because CNAs are inferred relative to the reference cells, we show how the choice of reference shapes the results: on healthy blood cells and normal tissue sections, false calls followed cell-type expression programmes that the reference did not cover, could be predicted from expression alone, and fell several-fold when the reference included the lineage of the cells analysed. We also show that it applies without change to long-read single-cell and spatial data: on cancer cell lines profiled with PacBio Kinnex, ONT and Illumina sequencing, its calls agreed with cell-line copy number to a similar degree for all three platforms, and on Kinnex Visium HD sections it mapped tumour subclones in tissue. inferCNV is available from Bioconductor, GitHub and Docker Hub.
